\section{Mid-Term Examination Suite (Units 1, 2 \& 3)}
\LPUPaperHeader{Mid-Term Examination}{Unified Model Paper}{90 Minutes}{30 Marks}{CO1 (DC), CO2 (AC), CO3 (Machines)}

\begin{center}
\sffamily\textbf{\color{AlertRed}Structure: Section A consists of 20 MCQs (1 Mark each). Section B consists of 2 Subjective Problems (5 Marks each).}
\end{center}

\subsection*{Section A: Multiple Choice Questions (20 Marks)}

\mcqitem{1}{The SI unit of electrical conductivity is:}{Ohm-meter}{Siemens/meter}{Henry/meter}{Farad/meter}{1}
\mcqitem{2}{Kirchhoff's Current Law is a mathematical consequence of the conservation of:}{Energy}{Linear Momentum}{Electric Charge}{Magnetic Flux}{1}
\mcqitem{3}{The Maximum Power Transfer theorem states that maximum power is transferred when load resistance equals:}{Zero}{Twice Thevenin resistance}{Thevenin internal resistance}{Infinity}{1}
\mcqitem{4}{Three equal resistors of $9\,\Omega$ are connected in delta. The resistance of each branch in the equivalent star network is:}{$27\,\Omega$}{$9\,\Omega$}{$3\,\Omega$}{$1\,\Omega$}{1}
\mcqitem{5}{The Form Factor of a pure sinusoidal waveform is approximately:}{$1.414$}{$1.11$}{$0.707$}{$0.637$}{1}
\mcqitem{6}{In a pure inductive AC circuit, the current:}{Leads voltage by $90^\circ$}{Lags voltage by $90^\circ$}{Is in phase with voltage}{Lags voltage by $180^\circ$}{1}
\mcqitem{7}{At series resonance in an $RLC$ circuit, the net reactance is:}{Inductive}{Capacitive}{Zero}{Infinite}{1}
\mcqitem{8}{In a balanced 3-phase delta-connected system, the line current $I_L$ relates to phase current $I_{ph}$ by:}{$I_L = I_{ph}$}{$I_L = \sqrt{3}I_{ph}$}{$I_L = I_{ph}/\sqrt{3}$}{$I_L = 3 I_{ph}$}{1}
\mcqitem{9}{If two wattmeters in a balanced three-phase measurement read identical positive values, the circuit power factor is:}{$0.5$ lagging}{Zero}{$1.0$ (Unity)}{$0.866$ leading}{1}
\mcqitem{10}{The primary winding and secondary winding of an ideal transformer are coupled:}{Electrically only}{Magnetically only}{Both electrically and magnetically}{Thermally only}{1}
\mcqitem{11}{The core of a transformer is laminated in order to significantly reduce:}{Hysteresis loss}{Copper loss}{Eddy current loss}{Stray load loss}{1}
\mcqitem{12}{If a transformer designed for $230\text{ V AC}$ is connected to a $230\text{ V DC}$ source:}{It works with high efficiency}{Secondary voltage doubles}{Primary draws excessive current and burns}{Core demagnetizes}{1}
\mcqitem{13}{In a DC motor, the back-EMF $E_b$ acts to:}{Assist the applied voltage}{Oppose the applied voltage}{Double the armature current}{Produce copper losses}{1}
\mcqitem{14}{A DC series motor must never be started without a mechanical load because:}{It will fail to start}{It will draw zero current}{Its speed will rise to dangerously high runaway levels}{Its back-EMF vanishes}{1}
\mcqitem{15}{Fleming's Left-Hand Rule is applied to determine the direction of force in:}{DC Generators}{DC Motors}{Transformers}{Capacitors}{1}
\mcqitem{16}{The speed of the rotating magnetic field in an induction motor is designated as:}{Rotor speed}{Slip speed}{Synchronous speed}{Base speed}{1}
\mcqitem{17}{At standstill ($N = 0$), the slip $s$ of an induction motor is:}{$0$}{$0.5$}{$1.0$}{$-1.0$}{1}
\mcqitem{18}{A single-phase induction motor is not inherently self-starting because:}{It lacks a rotor}{Its pulsating field produces zero net starting torque}{Its resistance is too high}{It has no stator}{1}
\mcqitem{19}{In an auto-transformer with turns ratio $k = V_2/V_1 = 0.8$, the copper saved compared to a two-winding transformer is:}{$20\%$}{$80\%$}{$50\%$}{$64\%$}{1}
\mcqitem{20}{An instrument transformer used to step down high line currents for standard meters is a:}{Potential Transformer}{Current Transformer}{Distribution Transformer}{Power Transformer}{1}

\newpage
\subsection*{Section B: Analytical Problems (10 Marks)}

\questionitem{21}{
    A $50\text{ kVA}$, $4400/220\text{ V}$, $50\text{ Hz}$ single-phase transformer has full-load copper loss of $600\text{ W}$ and iron loss of $450\text{ W}$. Calculate: (a) Efficiency at full load with $0.8$ power factor lagging, and (b) The kVA load at which maximum efficiency occurs, and the value of maximum efficiency at unity power factor.
}{5}{CO3}{L3 Apply}

\questionitem{22}{
    A $230\text{ V}$ DC shunt motor has an armature resistance $R_a = 0.5\,\Omega$ and field resistance $R_{sh} = 115\,\Omega$. When running on no load, it takes a total line current of $4\text{ A}$ at $1200\text{ rpm}$. Calculate its speed when running on full load drawing a line current of $30\text{ A}$, assuming the magnetic flux remains constant.
}{5}{CO3}{L3 Apply}

\newpage
\subsection*{Solutions and Marking Rubrics --- Mid-Term Examination}

\begin{solutionbox}[Answer Key: Section A (MCQs)]
\begin{center}\small
\begin{tabularx}{\linewidth}{c c c c c c c c c c}
\toprule
\textbf{Q1} & \textbf{Q2} & \textbf{Q3} & \textbf{Q4} & \textbf{Q5} & \textbf{Q6} & \textbf{Q7} & \textbf{Q8} & \textbf{Q9} & \textbf{Q10} \\
B & C & C & C & B & B & C & B & C & B \\
\midrule
\textbf{Q11} & \textbf{Q12} & \textbf{Q13} & \textbf{Q14} & \textbf{Q15} & \textbf{Q16} & \textbf{Q17} & \textbf{Q18} & \textbf{Q19} & \textbf{Q20} \\
C & C & B & C & B & C & C & B & B & B \\
\bottomrule
\end{tabularx}
\end{center}
\end{solutionbox}

\begin{solutionbox}[Solution to Question 21]
\textbf{Step 1: Efficiency at Full Load ($0.8$ pf)}\\
Full load active power: $P_{out} = 50 \times 0.8 = 40\text{ kW} = 40000\text{ W}$.\\
Total losses: $P_{loss} = P_i + P_{cu} = 450 + 600 = 1050\text{ W}$.\\
\[ \eta_{FL} = \frac{40000}{40000 + 1050} \times 100\% = \frac{40000}{41050} \times 100\% \approx 97.44\% \]
\textbf{Step 2: Maximum Efficiency Condition}\\
Load fraction for maximum efficiency: $x = \sqrt{P_i / P_{cu}} = \sqrt{450 / 600} = \sqrt{0.75} \approx 0.866$.\\
Load kVA at $\eta_{max}$: $S_{max} = 0.866 \times 50\text{ kVA} = 43.30\text{ kVA}$.\\
At unity power factor ($\cos\phi = 1.0$), output power is $P = 43300\text{ W}$.\\
Total loss at maximum efficiency $= 2 \times P_i = 2 \times 450 = 900\text{ W}$.\\
\[ \eta_{max} = \frac{43300}{43300 + 900} \times 100\% = \frac{43300}{44200} \times 100\% \approx 97.96\% \]
\begin{markingrubric}
$\bullet$ Full load efficiency computation: \textbf{2.5 Marks}\\
$\bullet$ Max efficiency fraction $x$, corresponding kVA, and $\eta_{max}$ value: \textbf{2.5 Marks}
\end{markingrubric}
\end{solutionbox}

\begin{solutionbox}[Solution to Question 22]
\textbf{Step 1: Field Current and No-Load Back EMF}\\
Shunt field current: $I_{sh} = \frac{V}{R_{sh}} = \frac{230}{115} = 2.0\text{ A}$.\\
No-load armature current: $I_{a0} = I_0 - I_{sh} = 4.0 - 2.0 = 2.0\text{ A}$.\\
No-load back EMF: $E_{b0} = V - I_{a0} R_a = 230 - (2.0 \times 0.5) = 230 - 1.0 = 229.0\text{ V}$.\\
\textbf{Step 2: Full-Load Armature Current and Back EMF}\\
Full-load line current: $I_{FL} = 30\text{ A} \implies I_{a1} = 30 - 2.0 = 28.0\text{ A}$.\\
Full-load back EMF: $E_{b1} = V - I_{a1} R_a = 230 - (28.0 \times 0.5) = 230 - 14.0 = 216.0\text{ V}$.\\
\textbf{Step 3: Speed Determination}\\
Since flux is constant ($\Phi_1 = \Phi_0$):
\[ \frac{N_1}{N_0} = \frac{E_{b1}}{E_{b0}} \implies N_1 = 1200 \times \frac{216.0}{229.0} \approx 1131.88\text{ rpm} \]
\begin{markingrubric}
$\bullet$ Correct calculation of $I_{sh}$ and armature currents: \textbf{2.0 Marks}\\
$\bullet$ Back-EMF calculations $E_{b0}$ and $E_{b1}$: \textbf{1.5 Marks}\\
$\bullet$ Ratio relation and final speed computation: \textbf{1.5 Marks}
\end{markingrubric}
\end{solutionbox}
\newpage
